% ============================================================
%  BioKGSuite — LaTeX Report (Nature-style compact)
% ============================================================
\documentclass[10pt,a4paper]{article}
% ---------- packages ----------
\usepackage[top=2.2cm,bottom=2.2cm,left=2cm,right=2cm]{geometry}
\usepackage{amsmath,amssymb}
\usepackage{graphicx}
\usepackage{booktabs}
\usepackage{array}
\usepackage{tabularx}
\usepackage{caption}
\usepackage{subcaption}
\usepackage{float}
\usepackage{setspace}
\usepackage{xcolor}
\usepackage{hyperref}
\usepackage[numbers,sort&compress]{natbib}
\usepackage{enumitem}
% ---------- Nature-style formatting ----------
\hypersetup{colorlinks=true,linkcolor=black,citecolor=blue!70!black,urlcolor=blue!70!black}
\setlength{\parskip}{4pt}
\setlength{\parindent}{0pt}
\captionsetup{font=small,labelfont=bf,skip=4pt}
% ---------- compact sections ----------
\usepackage{titlesec}
\titleformat{\section}{\normalsize\bfseries}{\thesection}{0.5em}{}
\titleformat{\subsection}{\small\bfseries}{\thesubsection}{0.5em}{}
\titlespacing*{\section}{0pt}{10pt}{4pt}
\titlespacing*{\subsection}{0pt}{8pt}{2pt}
% ============================================================
\begin{document}
\begin{center}
  {\large\textbf{BioKGSuite: A Multidimensional Framework \\ [0.25em]
  for the Evaluation of Biomedical Knowledge Graphs}}\\[0.5em]
  {\normalsize Emily Molins}\\[0.15em]
  {\normalsize March 2026}
\end{center}
\vspace{0.4em}
% ============================================================
\section{Introduction}
Biomedical knowledge graphs (KGs) integrate heterogeneous biological entities---drugs, genes, diseases, and pathways---into unified graph structures that support applications from drug repurposing to disease--gene prioritisation\citep{Yue2020, Su2023}. Over 120 are now catalogued\citep{KGHub2025}, yet their proliferation has outpaced rigorous quality assessment.

We present BioKGSuite, a reproducible framework that defines KG quality across seven dimensions---coverage, annotation accuracy, trustworthiness, topology, stability, task performance, and generalisation---operationalised via 19 metrics (Fig.~\ref{fig:framework}).

% ============================================================
\section{Framework}
BioKGSuite distinguishes three orthogonal aspects of KG quality: what information the graph contains (\textit{content}), how that information is organised (\textit{structure}), and how well it supports downstream biological inference (\textit{inference}).

\begin{figure}[H]
  \centering
  \includegraphics[width=\linewidth]{Figures/taxonomy.png}
  \caption{BioKGSuite evaluation framework: seven dimensions, 19 metrics, three categories.}
  \label{fig:framework}
\end{figure}

\textit{Content} dimensions assess what a KG contains and how well it is documented: \textbf{coverage} quantifies entity and relation breadth against gold-standard references; \textbf{annotation accuracy} evaluates identifier hygiene and schema conformance; and \textbf{trustworthiness} captures provenance and evidential support.

\textit{Structure} dimensions characterise how the graph is organised. \textbf{Topology} examines how well-connected and internally organised the graph is, assessing the extent to which nodes form a coherent cluster, remain mutually reachable, and group into communities that reflect meaningful biological relationships. \textbf{Stability} tests whether this topological signal is robust to missing data.

\textit{Inference} dimensions evaluate whether the graph supports accurate biological predictions. \textbf{Task performance} benchmarks downstream tasks including link prediction and multi-hop reasoning. \textbf{Generalisation} tests whether task performance holds across contexts.

Each dimension is scored on a scale from 0 to 1, where higher values indicate better performance. To demonstrate BioKGSuite, we evaluated four biomedical KGs---Hetionet, DRKG, PrimeKG, and OpenBioLink---with gold standards tailored to drug repurposing.

% ============================================================
\section{Results}

\begin{figure}[H]
  \centering
  \includegraphics[width=0.6\linewidth]{Figures/radar.png}
  \caption{Multi-dimensional quality profiles of the four evaluated KGs. No KG dominates all dimensions; each exhibits a distinct quality profile.}
  \label{fig:radar}
\end{figure}

\subsection{Content quality}
Annotation accuracy is near-perfect across all four KGs, suggesting that identifier hygiene and schema conformance have become baseline expectations rather than differentiating features.

The substantive variation lies in trustworthiness and coverage. Within trustworthiness, provenance granularity ranges from per-edge source attribution (DRKG, OpenBioLink) through per-node provenance (PrimeKG) to type-level metadata only (Hetionet); notably, no KG provides per-edge confidence scores, leaving uncertainty quantification unaddressed across the board. Coverage follows a similar but non-identical ordering, with PrimeKG the most comprehensive and Hetionet the most limited. That these two hierarchies diverge---PrimeKG leads on coverage but not provenance, while DRKG and OpenBioLink achieve the finest-grained attribution from a narrower entity base---indicates that breadth and traceability are independent dimensions of content quality and should be assessed separately.

\subsection{Structural properties}
All four KGs exhibit robust small-world topology, with clustering coefficients far exceeding random-graph baselines and near-complete node connectivity. Despite this shared foundation, meaningful structural differences emerge. OpenBioLink, despite being the largest KG by node count, has the lowest global reachability---a substantial fraction of its node pairs are not connected by any path, suggesting that large-scale heterogeneous integration can fragment global connectivity. Community structure reveals a related pattern: PrimeKG, organised around disease nodes, produces communities that align most closely with biological entity types, while DRKG, despite being purpose-built for drug repurposing, shows the weakest alignment---likely a consequence of its unusually broad relational schema.

Stability analysis corroborates these structural differences. All KGs are robust to random edge removal, confirming that predictive signal is broadly distributed rather than concentrated in a small number of critical connections. Periphery-targeted dropout, however, reveals sharper differentiation: Hetionet maintains near-identical performance under both strategies, reflecting the resilience of a compact and tightly curated graph, while OpenBioLink degrades substantially---consistent with its fragmented reachability profile and suggesting a structural dependence on hub connectivity.

\subsection{Inference capacity}
Across three biologically grounded tasks---link prediction of drug--disease indications, neighbourhood retrieval of disease-associated genes, and multi-hop reasoning over mechanistic drug--gene--disease paths---no single KG dominates. PrimeKG achieves the strongest overall profile, leading on both link prediction (AUROC 0.947) and multi-hop reasoning, while OpenBioLink achieves the highest neighbourhood retrieval. This task-dependent variation suggests that different KGs are structurally suited to different inference types, and that aggregate scores mask meaningful differences in predictive capability.

Hetionet presents the clearest case of coverage as a ceiling on inference. Its link prediction AUROC is competitive, but neighbourhood retrieval and multi-hop reasoning are the weakest of all four KGs by a substantial margin, consistent with a graph too narrow to support retrieval across the full diversity of disease--gene associations. Link prediction alone---the most commonly reported benchmark in the literature---overstates the inference capacity of coverage-limited KGs.

\subsection{Generalisation}
Generalisation performance is largely decoupled from task performance, though cross-KG comparisons should be interpreted carefully. Each KG is evaluated on the diseases, therapeutic areas, and entity pairs it actually contains, meaning differences in results partly reflect differences in what is being tested rather than performance alone.

Within these constraints, the data-sparse results are the most informative. DRKG shows the steepest degradation moving from well-annotated to low-connectivity diseases (40\% drop), indicating that its signal is disproportionately anchored to prevalent diseases. Hetionet degrades the least (8.2\%), though its disease graph is narrower and more evenly connected, which may itself reduce the severity of the sparse-entity setting. Cross-domain consistency follows a different ordering: DRKG and PrimeKG maintain stable performance across therapeutic areas (CV $\leq$ 1.9\%), while OpenBioLink shows nearly five-fold greater variance (CV = 9.6\%), suggesting its predictive signal is unevenly distributed across domains.

Prospective generalisation is evaluated only for KGs with $\geq$20 post-cutoff pairs, excluding Hetionet and OpenBioLink whose entity coverage is insufficient for a reliable estimate. Of the remaining two, DRKG performs no worse on prospective pairs than on contemporaneous ones, suggesting its relational breadth encodes structural signal for emerging indications rather than merely reflecting known associations. PrimeKG shows a modest but consistent decline, pointing to a degree of recency bias in its coverage.

\medskip
Figure~\ref{fig:heatmap} summarises sub-metric scores across all 19 metrics, revealing the full heterogeneity of quality profiles across KGs.

\begin{figure}[H]
  \centering
  \includegraphics[width=\linewidth]{Figures/heatmap.png}
  \caption{Sub-metric scores across all 19 BioKGSuite metrics. Darker cells indicate higher scores (0--1 scale).}
  \label{fig:heatmap}
\end{figure}

% ============================================================
\section{Discussion}
Taken together, these results reveal that biomedical KG quality is genuinely multidimensional, and that the dimensions are largely independent: strong performance on one axis does not predict performance on another. The most striking illustration is the dissociation between structural and inferential quality---all four KGs exhibit strong topological properties (small-world structure, near-complete connectivity, robust stability), yet task performance varies by more than 50\% across KGs. Graph structure alone does not determine inference capacity; entity coverage and relational richness set the effective ceiling.

Within the content category, the sharpest trade-offs lie between coverage and provenance. No KG achieves both broad entity coverage and fine-grained edge-level traceability, a tension that reflects the fundamental difficulty of large-scale heterogeneous integration. These are independent axes, and selecting a KG on coverage alone risks overlooking critical limitations in evidential support. The universal absence of per-edge confidence scores compounds this, and without uncertainty quantification, downstream models treat all edges as equally reliable.

These findings point to a practical conclusion. The right KG is application-dependent. Broad discovery campaigns may favour coverage; projects requiring auditable evidence chains should prioritise traceability; applications targeting rare diseases demand scrutiny of generalisation profiles. BioKGSuite is designed as a diagnostic instrument for exactly these decisions, a reproducible baseline against which future KG releases can be evaluated.

% ============================================================
\bibliographystyle{unsrtnat}
\begin{thebibliography}{9}

\bibitem[Caufield et al.(2023)]{KGHub2025}
Caufield, J. H. et al. (2023).
KG-Hub---building and exchanging biological knowledge graphs.
\textit{Bioinformatics}, 39(7):btad418.

\bibitem[Su et al.(2023)]{Su2023}
Su, C. et al. (2023).
Biomedical discovery through the integrative biomedical knowledge hub (iBKH).
\textit{iScience}, 6(4):106460.

\bibitem[Yue et al.(2020)]{Yue2020}
Yue, X. et al. (2020).
Graph embedding on biomedical networks: methods, applications and evaluations.
\textit{Bioinformatics}, 36(4):1241--1251.

\end{thebibliography}
\end{document}
